\documentclass{article}

% ---- graphiques
\usepackage[pdftex]{graphicx}
\usepackage{wrapfig}
\usepackage{color}
%\usepackage{hyperref}

% for latex2html
\usepackage{html}

% for accents
\usepackage[latin1]{inputenc}
\usepackage[T1]{fontenc}

\usepackage{algorithm}
\usepackage{algorithmic}

\definecolor{darkgreen}{rgb}{0,0.4,0}
\definecolor{darkblue}{rgb}{0,0,0.4}
\definecolor{darkgray}{rgb}{0.2,0.2,0.2}

% ---- inclusion de codes
\usepackage{listings}
\lstset{showstringspaces=false,tabsize=4,basicstyle=\scriptsize\sffamily,breaklines=true,breakatwhitespace=true,framexleftmargin=5mm, frame=shadowbox, framesep=1pt,rulesepcolor=\color{darkgray},rulesep=.5pt,keywordstyle=\bf\color{blue},commentstyle=\color{magenta},stringstyle=\color{red},numbers=left,numberstyle=\tiny,numbersep=5pt,columns=flexible}

\lstdefinestyle{bash}{language=bash}
\lstdefinestyle{Perl}{language=Perl}
\lstdefinestyle{C++}{language=C++,emph={__global__,__shared__,__syncthreads,blockIdx,threadIdx,float3,float4},emphstyle=\bf\color{darkgreen}}
\lstdefinestyle{DTD}{language=XML}
\lstdefinestyle{XML}{language=XML,usekeywordsintag=false,markfirstintag=true}
%begin{latexonly}
\newcommand{\includecode}[2]{
\lstinputlisting[style=#1]{#2}
}
%end{latexonly}
\begin{htmlonly}
\newcommand{\includecode}[2]{  \htmladdnormallink{#2}{../../#2} }
\end{htmlonly}

%\lstnewenvironment{code}{}{}
\lstnewenvironment{code_bash}{\lstset{style=bash}}{}
\lstnewenvironment{code_perl}{\lstset{style=Perl}}{}
\lstnewenvironment{code_cpp}{\lstset{style=C++}}{}
\lstnewenvironment{code_dtd}{\lstset{style=DTD}}{}
\lstnewenvironment{code_xml}{\lstset{style=XML}}{}

\newcommand{\textcode}[1]{{\sf #1}}



%
\newcommand{\sofa}{SOFA}
\newcommand{\todo}[1]{}
\newcommand{\eg}{\textit{e.g.} }

\renewcommand{\vec}[1]{\ensuremath{\mathbf{#1 }}} % vector
\newcommand{\Vx}{\vec{x} } % position vector
\newcommand{\Vv}{\vec{v} } % velocity vector
\newcommand{\Va}{\vec{a} } % acceleration vector
\newcommand{\Vf}{\vec{f}} % force
\newcommand{\Vdv}{\vec{\delta\Vv}} % change of velocity vector (unknown in implicit CG, and used in constraint solver
\renewcommand{\P}{\mat{P} } % projection to a constrained space.

\newcommand{\JNL}{\mathbf{\mathcal{J}} }     % mapping des positions
\newcommand{\J}{\mat J }                 % mapping lineaire
\newcommand{\M}{\mat M }             % matrice de masse
\newcommand{\K}{\mat K }             % matrice de raideur
\newcommand{\B}{\mat B }             % matrice d'amortissement
\newcommand{\G}{\mat G }             % jacobien des contraintes



% ---- inclusion de codes
\definecolor{darkgreen}{rgb}{0,0.4,0}
\definecolor{darkblue}{rgb}{0,0,0.4}
\definecolor{darkgray}{rgb}{0.2,0.2,0.2}


% macros mathematiques
\newcommand{\ma}[1]{\ensuremath{\mathbf {#1}}}
\newcommand{\ve}[1]{\ensuremath{\mathbf {#1}}}

\usepackage{amsmath}
\usepackage{amsfonts}
\usepackage{amssymb}

% character styles
\newcommand{\bm}[1]{\ensuremath{\mathbf{{#1}}}}
\newcommand{\mcal}[1]{\mbox{$\mathcal #1$}} % rondes math
\newcommand{\bmcal}[1]{\mbox{\boldmath $\mathcal #1$}} % rondes grasses math
\newcommand{\ensemble}[1]{\mbox{$\mathbb{#1}$}}
\newcommand{\RRR}{\mbox{$\ensemble{R}^3$}} 


% d�finitions
\newcommand{\definition}[2]{\index{#1}{\bf #1}: #2}
\newcommand{\voc}[1]{\index{#1}#1}
\newcommand{\bvoc}[1]{\index{#1}{\bf #1}}

% misc
\newcommand{\EV}[1]{\stackrel{\rightarrow}{#1}}  % espace vectoriel
\newcommand{\EA}[1]{#1}                          % espace affine

% vectors, matrices
%\newcommand{\point}[1]{\mbox{$#1$}}          % un point
\newcommand{\point}[1]{\ensuremath{#1}}          % un point
\newcommand{\mat}[1]{\bm{#1}}         % matrice
\newcommand{\matnm}[3]{\bm{#1_{#2\times #3}}}  % matrice n lignes , m colonnes
\newcommand{\vect}[1]{\bm{#1}}        % vecteur 
%\newcommand{\vecf}[1]{\stackrel{\rightarrow}{#1}}  % vecteur avec fleche
\newcommand{\vecf}[1]{\mbox{$\overrightarrow{#1}$}}  % vecteur avec fleche
\newcommand{\ident}[1]{\bm{I_{#1}}}   % identit� en dimension n
\newcommand{\inv}[1]{#1^{-1}}         % matrice inverse
\newcommand{\psinv}[1]{#1^{+}}        % matrice pseudo-inverse
\newcommand{\transp}[1]{#1^T}         % transpos�e de 1
\newcommand{\trace}[1]{tr(#1)}        % trace
\newcommand{\deter}[1]{\mbox{$|#1|$}}       % determinant
\newcommand{\oppvec}[1]{\mbox{$\left( \vect {#1} \wedge \right)$}}  % operateur matriciel de produit vectoriel

% bases, reperes
\newcommand{\vecin}[2]{\mbox{${}^{#2}#1$}}    % vecteur 1 dans repere 2
\newcommand{\Base}[1]{\ensuremath{\mathcal B_{#1}}} % Symbole du repere 1
\newcommand{\chbase}[3]{\mbox{${}_{#2}^{#3}\mat{#1}$}}  % operateur 1 fait le passage de la base 3 vers la base 2
%\newcommand{\pchbase}[2]{\chbase{\mat{B}}{#1}{#2}}  % matrice de passage de la base 2 vers la base 1
\newcommand{\pchbase}[2]{\chbase{B}{#1}{#2}}  % matrice de passage de la base 2 vers la base 1
\newcommand{\Rep}[1]{\ensuremath{\mathcal R_{#1}}} % Symbole du repere 1
\newcommand{\rep}[1]{\Rep{#1}}                 % Symbole du repere 1
%\newcommand{\pchrep}[2]{\chbase{\mat{F}}{#1}{#2}}  % matrice de passage du repere 1 vers le repere 2, F comme Frame
\newcommand{\pchrep}[2]{\chbase{\bm{C}}{#1}{#2}}  % matrice de passage du repere 2 vers le repere 1

%% Operateur de passage du repere 1 par rapport a 2
%\newcommand{\ChgRep}[2]{\mbox{\boldmath $R_{#1}^{#2}$}}

% rotations	
%\newcommand{\rot}[2]{\mbox{$\mat{R}_{#1,#2}$}}      % rotation vectorielle
\newcommand{\rot}[2]{\ensuremath{\mat{R}_{#1,#2}}}      % rotation vectorielle
\newcommand{\rota}[3]{\mbox{$\mat{R}_{#1,#2,#3}$}}  % rotation affine

% translation
\newcommand{\trans}[2]{\mbox{$\chbase{\vect{t}}{#1}{#2}$}} % passage de #1 vers #2 par une translation, ou translation du repere #2 par rapport au repere #1

% vitesses et acc�l�rations
\newcommand{\VRep}[2]{\mbox{\boldmath $\dot R_{#1}^{#2}$}} % vitesse du repere 1 par rapport a 2 
%\newcommand{\Point}[2]{\mbox{\boldmath ${#1}^{#2}$}}  % Coordonnees d'un point 1 dans un repere 2
\newcommand{\Point}[2]{\mbox{$\vecin{\bm{#1}}{#2}$}}  % Coordonnees d'un point 1 dans un repere 2
\newcommand{\VPoint}[2]{\mbox{\boldmath ${\dot #1}_{/#2}$}} % Vitesse d'un point par rapport � un repere
\newcommand{\APoint}[2]{\mbox{\boldmath ${\ddot #1}_{/#2}$}} % Acceleration d'un point par rapport � un repere

% cinematique du solide
\newcommand{\derivedans}[2]{\mbox{$\dot{#1}^{(#2)}$}}  % derivee du vecteur 1 dans repere 2
\newcommand{\fixedans}[2]{\mbox{$#1_{\in #2}$}}        % vecteur 1 fixe dans repere 2
\newcommand{\vecom}{\mbox{$\bm{\Omega}$}}  % omega de 1 par rapport a 2
\newcommand{\vecrot}[2]{\mbox{$\vecom_{#1/#2}$}}  % omega de 1 par rapport a 2
\newcommand{\accrot}[2]{\mbox{$\dot{\vecom}_{#1/#2}$}}  % omega de 1 par rapport a 2
\newcommand{\vfdans}[3]{\mbox{$\vec V^{#2/#3}_{#1}$}}    % vitesse de 1 fixe dans 2 par rapport a 3
\newcommand{\afdans}[3]{\mbox{$\vec \Gamma^{#2/#3}_{#1}$}}    % acceleration de 1 fixe dans 2 par rapport a 3
\newcommand{\vmdans}[2]{\mbox{$\vec V^{/{#2}}_{#1}$}}    % vitesse de 1 mobile dans 2
\newcommand{\amdans}[2]{\mbox{$\vec \Gamma^{/#2}_{#1}$}}    % acceleration de 1 mobile dans 2

% chaines articulees
\newcommand{\liaison}[2]{\mbox{$\mathcal L_{#1,#2}$}}  % liaison du pere 1 vers fils 2 (et repere intermediaire)
\newcommand{\liaisonprime}[2]{\mbox{$\mathcal L'_{#1,#2}$}}  % deuxieme repere intermediaire de la liaison du pere 1 vers fils 2
\newcommand{\liaisonP}[2]{\mbox{$\mathcal L_{#1,#2}$}}  % Repere dans pere 1 de la liaison vers fils 2 
\newcommand{\liaisonC}[2]{\mbox{$\mathcal L'_{#1,#2}$}}  % Repere dans fils de la liaison du pere 1 vers fils 2 
%\newcommand{\transP}[2]{\pchrep{\liaisonP{#1}{#2}}{#1}}  % Matrice du repere dans pere de la liaison du pere 1 vers fils 2 
%\newcommand{\transC}[2]{\pchrep{\liaisonC{#1}{#2}}{#2}}  % Matrice du repere dans pere de la liaison du pere 1 vers fils 2 
%\newcommand{\transPC}[2]{\pchrep{\liaisonC{#1}{#2}}{\liaisonP{#1}{#2}}}  % matrice de passage entre repere liaison dans fils et repere de liaison dans pere
\newcommand{\transP}[2]{\chbase{C_p}{#2}{#1}}  % Matrice du repere dans pere de la liaison du pere 1 vers fils 2 
\newcommand{\transC}[2]{\chbase{C_c}{#2}{#1}}  % Matrice du repere dans pere de la liaison du pere 1 vers fils 2 
\newcommand{\transPC}[2]{\chbase{C_l}{#2}{#1}}  % matrice de passage entre repere liaison dans fils et repere de liaison dans pere
% \pchrep{fils}{pere} = \liaisonP{pere}{fils}\deplPC{pere}{fils}\liaisonC{pere}{fils}

%topology
\newcommand{\mesh}{{\mathcal M}}
\newcommand{\vertices}{{\mathcal V}}
\newcommand{\edges}{{\mathcal E}}
\newcommand{\triangles}{{\mathcal TR}}
\newcommand{\tetrahedra}{{\mathcal T}}
\newcommand{\controls}{{\mathcal C}}
\newcommand{\nvertices}{{ V}}
\newcommand{\nedges}{{ E}}
\newcommand{\ntriangles}{{ TR}}
\newcommand{\ntetrahedra}{{ TE}}
\newcommand{\ncontrols}{{C}}
\newcommand{\control}{{\mathbf C}}
\newcommand{\degree}{{d}}
\newcommand{\euc}{{\rm I\!R}}
\newcommand{\naturalSet}{{\rm I\!N}}

\newcommand{\pctab  }{\hspace{0.15in}      }  % Pseudo-code indentation.
\newcommand{\code}[1]{ 
\begin{makeimage}
\begin{tabbing} \pctab \= \pctab \= \pctab \= \pctab \= \pctab \= \pctab \= \pctab \kill
#1
\end{tabbing}
\end{makeimage}
}



\newenvironment{componentoption}[1]%
{\textbf{#1}\newline}
{\newline}

\newcommand{\aliases}[1] {\newline \textit{Aliases - } #1}
%\newcommand{\defaultvalue}[1] {\newline \textit{Default Value - } #1}
\newcommand{\valuetype}[1] {\newline \textit{Value Type - } \textbf{#1}}

\usepackage{graphicx}
\usepackage{qtree}

\begin{document}
\raggedright

\title{SofaPython plugin}
\author{Bruno Carrez}



\maketitle

\begin{abstract}
The goal of this plugin is to provide python integration into \sofa. 

Therefore, it is possible to add new features in \sofa without modifying \sofa codebase. Moreover, python scripts are stored in the same place as the corresponding scene, for a cleaner integration of scene-specific code without polluting the rest of \sofa codebase.

This plugin provides a \textcode{Sofa} module available from python scripts, to interact with \sofa. 
\end{abstract}

\begin{center}
\includegraphics{SofaPython.png}
\end{center}

%\
\section{Requirements}

SOFA Packages: 
The following must be enabled in sofa-local.prf
\begin{itemize}
\item Any special things that need to be enabled
\end{itemize}

SOFA Plugins:
The following must be loaded in your SOFA instance
\begin{itemize}
\item Any plugins that must be loaded
\end{itemize}

\section{Scene Settings}

\subsection{Required Settings}

\begin{componentoption}{RequiredExampleSetting}
A description of what this setting controls, and any special information the user should know about it. Required settings are those that need to be specified by the user in order for the component to function. If your component doesn't have any required settings, leave this section blank. The below "value type" is the type that is expected by the component. "Aliases" are other strings that can be specified in the scene file for this setting. These are defined in the code using "addAlias()".
\valuetype{string}
\aliases{requiredexamplesetting}
\end{componentoption}


\subsection{Optional Settings}

\begin{componentoption}{OptionalExampleSetting}
A description of what this setting controls, and any special information the user should know about it. Optional settings can be specified by the user in order to change the default behaviour of the component. The below "Default Value" is the value given to the setting if the user doesn't specify anything.
\valuetype{bool}
\defaultvalue{false}
\aliases{optionalexample, optionalsetting}
\end{componentoption}

\section{Scene Data}

\subsection{Required Data}

\begin{componentoption}{RequiredExampleData}
Data is usually a link to something in another component. Required data must be specified in order for the component to function.
\valuetype{ExampleType}
\end{componentoption}

\subsection{Optional Data}

\begin{componentoption}{OptionalExampleData}
Data is usually a link to something in another component.
\valuetype{ExampleType}
\aliases{OptData}
\end{componentoption}

\subsection{Output data}
Output data is generally not defined in the component, but is linked to by other components.

\subsection{Example File}
path/to/an/example/file.scn



\section{"PythonScriptController" component}




\
\section{SofaPython graph components}

\subsection{PythonScriptController}
This controller is an empty Controller, its behavior is driven in the Python script it loads. THis is the simplest and the most polyvalent way to imbed Python scripting in a \sofa scene graph. 

   

\
\section{SofaPython scene files}

It is possible, once the plugin is charged in runSofa, to load python script *.py files directly instead of xml *.scn files. This way, you can create scene graphs in a procedural way instead of a descriptive way.

An example of this is provided in the \textcode{examples/PythonScene.py} sample scene of the plugin. 
   

\
\section{Creating graph objects in Python}


Creating objects in Python is as simple as describing them in XML scene files.

You can use either \textcode{Sofa.createObject(typename,...)} or \textcode{Sofa.BaseContext.createObject(typename,...)} to do it (the second creates the object THEN adds it to the context, whereas the first creates an "orphan" object taht you must add later to a context.)

Both functions take a string for the type name of the object to create, then any optional named parameters neede to complete the object description.

Here is an example of xml file, and its equivament in SofaPython code: (xml is commented)

\begin{code_python}
#        <EulerImplicit name="cg_odesolver" printLog="false" />
	node.createObject('EulerImplicit',name='cg_odesolver',printLog='false')
#        <CGLinearSolver iterations="25" name="linear solver" tolerance="1.0e-9" threshold="1.0e-9" />
	node.createObject('CGLinearSolver',name='linear solver',iterations=25,tolerance=2.0e-9,threshold=1.0e-9)
#        <MechanicalObject />
	node.createObject('MechanicalObject')
#        <UniformMass totalmass="10" />
	node.createObject('UniformMass',name='mass',totalmass=10)
#        <RegularGrid nx="6" ny="5" nz="3" xmin="-11" xmax="11" ymin="-7" ymax="7" zmin="-4" zmax="4" />
	node.createObject('RegularGrid', name='RegularGrid1', nx=6, ny=5, nz=3, xmin=-11, xmax=11, ymin=-7, ymax=7, zmin=-4, zmax=4)
#        <RegularGridSpringForceField name="Springs" stiffness="350" damping="1" />
	node.createObject('RegularGridSpringForceField',name='Springs',stiffness='350', damping='1')

#        <Node name="VisuDragon" tags="Visual">
	VisuNode = node.createChild('VisuDragon')
#            <OglModel name="Visual" filename="mesh/dragon.obj" color="red" />
	VisuNode.createObject('OglModel',name='Visual',filename='mesh/dragon.obj', color=color)
#            <BarycentricMapping object1="../.." object2="Visual" />
	VisuNode.createObject('BarycentricMapping',object1='../..', object2='Visual' )
#        </Node>	

#        <Node name="Surf">
	SurfNode = node.createChild('Surf')
#            <MeshObjLoader name="loader" filename="mesh/dragon.obj" />
	SurfNode.createObject('MeshObjLoader', name='loader', filename='mesh/dragon.obj')
#            <Mesh src="@loader" />
	SurfNode.createObject('Mesh',src='@loader')
#            <MechanicalObject src="@loader" />
	SurfNode.createObject('MechanicalObject', name='meca', src='@loader')
#            <Triangle />
	SurfNode.createObject('Triangle')
#            <Line />
	SurfNode.createObject('Line')
#            <Point />
	SurfNode.createObject('Point')
#            <BarycentricMapping />
	SurfNode.createObject('BarycentricMapping')
#        </Node>	


\end{code_python}


Remark: the values given for each optional named parameter can be a string, a scalar or an integer; this way you can flawlessly use computed variables, or even method calls to initialize your components.

\

\section{\sofa \ Python API}

\subsection{The \textcode{Sofa} module}

The core of this plugin is the \textcode{Sofa} Python module available to python scripts from within the SofaPython components (they are not available outside Sofa environment, in the command-line \textcode{python} binary for example).

Therefore, each python script to be imbedded in \sofa should include the following line if it wants to interact with the \sofa framework:

\begin{code_python}
import Sofa
\end{code_python}

This module provides a wide range of methods and types, bound to essential \sofa \ framework features.

These can be used from python within scripts loaded by the components provided by this plugin.

%Module methods 
\subsection{Module methods}

SofaPython provides several module methods, for general purpose (not linked to a particular node or component, for example).

\begin{itemize}
\item \textcode{sendGUIMessage(msgType,msgValue)} 
\item \textcode{createObject(type[, datafield=value [,...] ])} creates an object. ex: \textcode{node.createObject('UniformMass',totalmass=1)} See "Creating graph objects in Python" section for more info.
\end{itemize}

%Class Hierarchy
\subsection{Types hierarchy}

The class hierarchy in the \textcode{Sofa} module is quite different from the C++ \sofa \ class hierarchy. Not all \sofa \ classes are bound in Python, and some levels in the hierarchy are skipped.

Despite the ability of Python to support multi-heritage, this feature has not been implemented in the \textcode{Sofa} module, for code simplicity.


\begin{figure}[htbp]
\begin{center}
\Tree [.Base [.BaseContext [.Context [.BaseNode Node ] ] ] 
			[.BaseObject [.BaseState [.BaseMechanicalState MechanicalObject ] ] [.BaseLoader MeshLoader ] 
			[.Topology [.BaseMeshTopology [.MeshTopology GridTopoligy ] ] ] ] ]
\Tree [.Vector3 ]
\Tree [.LinearSpring ]
\Tree [.Data DisplayFlagsData ]
\caption{Sofa python types hierarchy}
\label{default}
\end{center}
\end{figure}

\newpage

%Sofa.Vector3
\subsubsection{Sofa.Vector3}
Attributes :
\begin{itemize}
\item \textcode{x} 
\item \textcode{y}
\item \textcode{z}
\end{itemize}

%Sofa.LinearSpring
\subsubsection{Sofa.LinearSpring}
Attributes :
\begin{itemize}
\item \textcode{Index1} : the first extremity of the spring
\item \textcode{Index2} : the second extremity of the spring
\item \textcode{Ks} : spring stiffness
\item \textcode{Kd} : damping factor
\item \textcode{L} : rest length of the spring
\end{itemize}

%Sofa.Topology
\subsubsection{Sofa.Topology}
Methods :
\begin{itemize}
\item \textcode{haspos()} 
\item \textcode{getNbPoints()}
\item \textcode{setNbPoints(n)}
\item \textcode{getPX(i)} 
\item \textcode{getPY(i)} 
\item \textcode{getPZ(i)} 
\end{itemize}

%Sofa.BaseMeshTopology
\subsubsection{Sofa.BaseMeshTopology}
Methods :
\begin{itemize}
\item \textcode{getNbEdges()} 
\item \textcode{getNbTriangles()}
\item \textcode{getNbQuads()}
\item \textcode{getNbTetrahedra()}
\item \textcode{getNbHexahedra()}
\end{itemize}

%Sofa.GridTopology
\subsubsection{Sofa.GridTopology}
Methods :
\begin{itemize}
\item \textcode{setSize(nx,ny,nz)} 
\item \textcode{setNumVertices(nx,ny,nz)} 
\item \textcode{getNx()}
\item \textcode{getNy()}
\item \textcode{getNz()}
\item \textcode{setNx(n)}
\item \textcode{setNy(n)}
\item \textcode{setNz(n)}
\end{itemize}

%Sofa.RegularGridTopology
\subsubsection{Sofa.RegularGridTopology}
Methods :
\begin{itemize}
\item \textcode{setPos(xmin,xmax,ymin,ymax,zmin,zmax))} 
\end{itemize}

%BASE
\subsubsection{Sofa.Base}
Methods :
\begin{itemize}
\item \textcode{findData(name)} returns a \textcode{Sofa.Data} object (if it exists)
\end{itemize}
Attributes:
\begin{itemize}
\item \textcode{name} 
\end{itemize}

\subsubsection{Sofa.BaseObject}
Methods :
\begin{itemize}
\item \textcode{init()}
\item \textcode{bwdInit()}
\item \textcode{reinit()}
\item \textcode{storeResetState()}
\item \textcode{reset()}
\item \textcode{cleanup()}
\item \textcode{getContext() returns a \textcode{Sofa.BaseContext} object}
\item \textcode{getMaster() returns a \textcode{Sofa.BaseObject} object}
\item \textcode{setSrc(valueString, loader)} use a \textcode{Sofa.BaseLoader} to initialize the object
\end{itemize}

\subsubsection{Sofa.BaseState}
Methods :
\begin{itemize}
\item \textcode{resize(size)}
\end{itemize}

\subsubsection{Sofa.BaseMechanicalState}
Methods :
\begin{itemize}
\item \textcode{applyTranslation(x,y,z)}
\item \textcode{applyRotation(x,y,z)} Euler angles in degrees.
\item \textcode{applyScale(x,y,z)} 
\end{itemize}

\subsubsection{Sofa.MechanicalObject}
Methods :
\begin{itemize}
\item \textcode{setTranslation(x,y,z)} Initial transformations accessor.
\item \textcode{setRotation(x,y,z)} Initial transformations accessor. Euler angles in degrees.
\item \textcode{setScale(x,y,z)} Initial transformations accessor.
\item \textcode{getTranslation()} returns a \textcode{Sofa.Vector3} object. Initial transformations accessor.
\item \textcode{getRotation()} returns a \textcode{Sofa.Vector3} object. Initial transformations accessor. Euler angles in degrees.
\item \textcode{getScale()} returns a \textcode{Sofa.Vector3} object. Initial transformations accessor.
\end{itemize}

\subsubsection{Sofa.BaseContext}
Methods :
\begin{itemize}
\item \textcode{getRootContext()} 
\item \textcode{getTime()} 
\item \textcode{getDt()} 
\item \textcode{getGravity()} returns a \textcode{Sofa.vec3} object
\item \textcode{setGravity(Vec3)} 
\item \textcode{createObject(type[, datafield=value [,...] ])} creates an object. ex: \textcode{node.createObject('UniformMass',totalmass=1)} See "Creating graph objects in Python" section for more info.
\item \textcode{getObject(path)}
\item \textcode{getObjects()}  returns the list of components in this context (locally) 
\end{itemize}
Attributes:
\begin{itemize}
\item \textcode{active}  (boolean)
\item \textcode{animate}  (boolean)
\end{itemize}


\subsubsection{Sofa.Context}
(empty)
\subsubsection{Sofa.Node}
\begin{itemize}
\item \textcode{createChild(childName)} 
\item \textcode{getRoot()} returns the root node of the graph (if any) 
\item \textcode{getChild(path)} returns a child node, given its path (if any) 
\item \textcode{getChildren()} returns the list of children 
\item \textcode{getParents()} returns the list of parents
\item \textcode{executeVisitor(visitor)} executes a python visitor from this node. See "Visitors" section for more infos.
\item \textcode{simulationStep(dt)} executes ONE step of simulation, using dt (in seconds, float) as delta time. 
\item \textcode{init()} Initialize the scene 
\item \textcode{reset()} Reset the scene 
\item \textcode{addObject(object)} 
\item \textcode{removeObject(object)} 
\item \textcode{addChild(childNode)} 
\item \textcode{removeChild(childNode)} 
\item \textcode{moveChild(childNode)} 
\item \textcode{detachFromGraph()} 
\item \textcode{sendScriptEvent(eventName, data)} propagates a script event from this node; eventName is the name of the event (string) and data is whatever you want: scalar, float, array, ... as long as it is a python object or structure. Other PythonScriptController will then receive the correponding \textcode{onScriptEvent(senderName, eventName, data)}. Useful for inter-script communications (see ScriptEvent.scn sample scene in the examples folder).
\end{itemize}


%DATA
\subsubsection{Sofa.Data}
Attributes :
\begin{itemize}
\item value
\item name
\end{itemize}
Methods :
\begin{itemize}
\item getValue(index)
\item setValue(index,value)
\end{itemize}

%DATA
\subsubsection{Sofa.DisplayFlagsData}

This class is always used in a "VisualStyle' object in the root node; see example below. It is used to control what should be displayed.

Attributes :
\begin{itemize}
\item showAll
\item showVisual
\item showVisualModels
\item showBehavior
\item showBehaviorModels
\item showForceFields
\item showInteractionForceFields
\item showCollision
\item showCollisionModels
\item showBoundingCollisionModels
\item showMapping
\item showMappings
\item showMechanicalMappings
\item showOptions
\item showWireFrame
\item showNormals
\end{itemize}
Unlike in C++ code, all these attributes are boolean values for simplicity.
Some attributes (like showAll) are hierarchicaly above others: setting them to \textcode{True} or \textcode{False} will set children attributes also; reading them will return \textcode{True} if their real C++ state is true or neutral.

\begin{figure}[htbp]
\begin{center}
\Tree [.All [.Visual VisualModels  ] 
			[ .Behavior BehaviorModels ForceFields Interactions ]
			[ .Collision CollisionModels BoundingCollisionModels ]
			[ .Mapping Mappings MechanicalMappings ]
			[ .Options Wireframe Normals ] ]
\caption{display flags hierarchy}
\label{default}
\end{center}
\end{figure}



Simple example of use:
\begin{code_python}
 style = node.getRoot().createObject('VisualStyle')
 style.findData('displayFlags').showBehaviorModels = True
\end{code_python}


\subsubsection{Data members: the most important thing in \sofa Python API}

The most important class is \textcode{Sofa.Base}, and its associated method \textcode{findData}.
ALMOST EVERYTHING in \sofa is stored in Datas, and with the only \textcode{Sofa.Base.findData} method, almost everything is possible.
Through the \textcode{Sofa.Data class} (returned by \textcode{Sofa.Base.findData(name)} ) it is possible to read or write almost any object value, thus interact with the simulation in real-time.

This way, even if a specific component isn't bound to python, it's possible to access it by its \textcode{Sofa.Base} heritage.

\textcode{Data.value} attribute has a versatile behavior, depending on the Data type.

On read, \textcode{Data.value} can return either an integer, a float, a string, or even a list of one of these 3 types.
On write, you have to set EXACTLY the proper type, or you can set a string (same format as in the *.scn xml files).

Examples:

\begin{code_python}
print str(node.findData('gravity').value
\end{code_python}
will output the text conversion of a list of 3 floats :
\begin{code_bash}
[0.0, -9.81, 0.0]
\end{code_bash}

You can set it in two ways ; the native version :
\begin{code_python}
node.findData('gravity').value = [0.0, -9.81, 0.0]
\end{code_python}
or by the text version:
\begin{code_python}
node.findData('gravity').value = '0.0  -9.81 0.0'
\end{code_python}

Use of any other type will result in an error.
The following won't work for example :
\begin{code_python}
node.findData('gravity').value = 9.81
\end{code_python}


\
\section{Visitors}

\subsection{Example visitor implementation in Python}

Following is a very simple visitor implementation in a python script:
(it is provided in the \textcode{examples/Visitor.scn} sample scene of the plugin)

\begin{code_python}
class SofaVisitor(object):
 def __init__(self,name):
  print 'SofaVisitor constructor name='+name
  self.name = name
  
 def processNodeTopDown(self,node):
  print 'SofaVisitor "'+self.name+'" processNodeTopDown node='+node.findData('name').value
  return True
 
 def processNodeBottomUp(self,node):
  print 'SofaVisitor "'+self.name+'" processNodeBottomUp node='+node.findData('name').value

\end{code_python}

Your visitor class MUST implement the two methods \textcode{processNodeTopDown} and \textcode{processNodeBottomUp}. 

In this example it is only meant to go through the graph in both directions, listing all node names.

The scene graph is the following nodes tree:

\begin{figure}[htbp]
\begin{center}
\Tree [.God [.Adam Abel  ] 
			Eve ]
\caption{Scene graph of the Visitor sample.}
\label{default}
\end{center}
\end{figure}

The Visitor is executed on the "God" node:
\begin{code_python}
 v = SofaVisitor('PythonVisitor')
 node.executeVisitor(v)
\end{code_python}

and the console output result is:

\textcode{
SofaVisitor constructor name=PythonVisitor
SofaVisitor "PythonVisitor" processNodeTopDown node=god
SofaVisitor "PythonVisitor" processNodeTopDown node=Adam
SofaVisitor "PythonVisitor" processNodeTopDown node=Abel
SofaVisitor "PythonVisitor" processNodeBottomUp node=Abel
SofaVisitor "PythonVisitor" processNodeBottomUp node=Adam
SofaVisitor "PythonVisitor" processNodeTopDown node=Eve
SofaVisitor "PythonVisitor" processNodeBottomUp node=Eve
SofaVisitor "PythonVisitor" processNodeBottomUp node=god}

\

\section{Sofa "PythonScriptController" component}

\subsection{Example scene}
\begin{code_bash}
runSofa applications/plugins/SofaPython/examples/ExampleController.scn
\end{code_bash}

\subsection{Component data}
\begin{itemize}
\item \textcode{filename} is the name of the .py script file.
\item \textcode{classname} is the name of the class in the script to instanciate for the controller. NOTE: the class MUST inherit from \textcode{Sofa.PythonScriptController}. See example files.
\end{itemize}

\subsection{Python script entry points}

\begin{itemize}
\item \textcode{onKeyPressed(self,c)}
\item \textcode{onKeyReleased(self,c)}
\item \textcode{onLoaded(self,node)}
\item \textcode{onMouseButtonLeft(self,mouseX,mouseY,isPressed)}
\item \textcode{onMouseButtonRight(self,mouseX,mouseY,isPressed)}
\item \textcode{onMouseButtonMiddle(self,mouseX,mouseY,isPressed)}
\item \textcode{onMouseWheel(self,mouseX,mouseY,wheelDelta)}
\item \textcode{onGUIEvent(self,strControlID,valueName,strValue)}
\item \textcode{onBeginAnimationStep(self,deltaTime)}
\item \textcode{onEndAnimationStep(self,deltaTime)}
\item \textcode{onScriptEvent(self,senderNode, eventName, data)}
\item \textcode{createGraph(self,node)}
\item \textcode{initGraph(self,node)}
\item \textcode{bwdInitGraph(self,node)}
\item \textcode{storeResetState(self)}
\item \textcode{reset(self)}
\item \textcode{cleanup(self)}
\end{itemize}





\end{document}
